% The first argument is both the visible label and the \Mystery{...} reference
% key. Each mystery therefore needs only a label and its body.
\MysterySheet{myst:how-was-newcomer-killed}{How was newcomer killed?}{

How was newcomer killed?

The answer is that the forbidden potion \Character{vampirehead} gave it killed it, due to \Character{usurper} having given it previously.
}

\MysterySheet{myst:who-killed-newcomer}{Who killed newcomer?}{

The newcomer was killed before he could be turned, who would to such a thing?

This mystery is supposed to be, in practice, who gave

The answer is \Character{usurper}.

}

\MysterySheet{myst:what-happend-to-newcomer}{What happened to the newcomer?}{

Did \Character{vampirehead} kill \Character{newcomer}?

Was the potion \Character{vampirehead} drink really the correct one?

The answer is that \Character{vampirehead} really only gave \Character{newcomer} the correct potion.

}

\MysterySheet{myst:what_happened_to_the_memories}{Why did all the vampires lose memories?}{

At some point in the game, the players are supposed to notice that all of them have missing memories.

}

\MysterySheet{myst:was_newcomer_of_family}{Was newcomer already a member of the family?}{

There are clues that this is possible and relevant.

}

% Read this as: solve the first mystery before the second mystery.
\MysteryDependency{myst:how-was-newcomer-killed}{myst:who-killed-newcomer}

\MysteryDependency{myst:what-happend-to-newcomer}{myst:how-was-newcomer-killed}
